\usetikzlibrary{arrows.meta,bbox,bending,calc,decorations.markings,decorations.pathreplacing,fpu,graphs,patterns.meta,quotes}
% utilities
% stolen from https://tex.stackexchange.com/a/463819
\newcommand{\tikz@fmgz@isnumber}[3]{% \tikz@fmgz@isnumber{<Number>}{<True>}{<False>}
    \begingroup
    \pgfkeys{/pgf/fpu/handlers/invalid number/.code={}}%
    \pgfmathfloatparsenumber{#1}%
    \global\pgfmathfloatgetflagstomacro\pgfmathresult\pgfFloatParseFlag%
    \endgroup
    \ifnum\pgfFloatParseFlag=3%
        #3%
    \else%
        #2%
    \fi%
}
% some private routines
\newcount\c@fmgz@counta
\newcount\c@fmgz@countb
\newif\if@tikz@fmgz@debug
\newcommand{\tikz@fmgz@typeout}[1]{\if@tikz@fmgz@debug
\typeout{#1}\fi}
\def\tikz@fmgz@pgfpoint#1,#2\pgf@stop{\pgfpoint{#1}{#2}}
\def\tikz@fmgz@pgfdist#1,#2;#3,#4\pgf@stop{\pgfmathparse{sqrt((#1-#3)*(#1-#3)+(#2-#4)*(#2-#4))}}
\def\tikz@fmgz@pgfangle#1,#2;#3,#4\pgf@stop{\pgfmathparse{atan2(#4-#2,#3-#1)}}
\newif\if@tikz@fmgz@firstup
\newif\iftikz@fmgz@paralell@curve@has@label
\tikzset{fmg/.style={/tikz/Feynmangraphz/install tikz keys,
	/tikz/graphs/use existing nodes,
	/tikz/Feynmangraphz/.cd,every diagram,#1,/tikz/.cd},
	Feynmangraphz/.cd,
		install tikz keys/.code={
			\tikzset{Gluon/.style={%/pgf/fpu/install only={reciprocal},
				draw,decorate,decoration={fmgz gluon,segment length=5pt,amplitude=6pt},propagator
				},
				gluon/.style={/tikz/draw=none,/tikz/decoration={name=none},postaction={Gluon}},
				Photon/.style={%/pgf/fpu/install only={reciprocal},
					draw,decorate,decoration={fmgz photon,segment length=5pt,amplitude=6pt,#1},propagator
				},
				Photon'/.style={%/pgf/fpu/install only={reciprocal},
					draw,decorate,decoration={fmgz photon',segment length=5pt,amplitude=6pt,#1},propagator
				},
				gauge boson/.style={/pgf/fpu/install only={reciprocal},/tikz/draw=none,/tikz/decoration={name=none},
					postaction={/tikz/Feynmangraphz/segment length=9pt,/tikz/Feynmangraphz/amplitude=1.5pt,wiggles,decorate}},
				photon/.style={/tikz/draw=none,/tikz/decoration={name=none},postaction={Photon={#1}}},
				photon'/.style={/tikz/draw=none,/tikz/decoration={name=none},postaction={Photon'={#1}}},
				vertex/.style={circle,inner sep=0.4ex,outer sep=0pt,fill,node contents={}},
				blob/.style={circle,inner sep=1ex,outer sep=0pt,draw,thick,pattern={Lines[angle=45,distance=2pt]},node contents={}},
				wiggles/.style={decoration={
						markings,
						mark=at position 0 with {\global\tikz@fmgz@n-1\relax
							\pgfmathtruncatemacro{\tikz@fmgz@k}{round((6*\pgfdecoratedpathlength+1*\pgfkeysvalueof{/tikz/Feynmangraphz/segment length})/(6*\pgfkeysvalueof{/tikz/Feynmangraphz/segment length}))}%	
							\pgfmathtruncatemacro{\tikz@fmgz@nn}{6*\tikz@fmgz@k}% number of steps
							\pgfmathsetmacro{\tikz@fmgz@step}{1/\tikz@fmgz@nn}%
							\xdef\tikz@fmgz@step{\tikz@fmgz@step}%
							%\pgfmathsetmacro{\tikz@fmgz@nn}{1/\tikz@fmgz@step}% just for debugging
							\tikz@fmgz@typeout{path length=\pgfdecoratedpathlength, k=\tikz@fmgz@k, step length=\tikz@fmgz@step, expected step number=\tikz@fmgz@nn}},
						mark=between positions 0 and 1 step \tikz@fmgz@step with {
							\global\advance\tikz@fmgz@n1\relax
							%\tikz@fmgz@typeout{\tikz@fmgz@tmp@node-m-\pgfkeysvalueof{/pgf/decoration/mark info/sequence number}}
							\path (0,{\pgfkeysvalueof{/tikz/Feynmangraphz/amplitude}*sqrt(abs(sin(\tikz@fmgz@n*60)))*sign(sin(\tikz@fmgz@n*60))})
								coordinate(\tikz@fmgz@tmp@node-m-\the\tikz@fmgz@n);
							},
						mark=at position 1 with {%\global\advance\tikz@fmgz@n1\relax
							%\path (0,0) coordinate(\tikz@fmgz@tmp@node-m-\the\tikz@fmgz@n);
							%\path (0,{\pgfkeysvalueof{/tikz/Feynmangraphz/amplitude}*sqrt(abs(sin(\tikz@fmgz@n*60)))*sign(sin(\tikz@fmgz@n*60))})
							%	coordinate(\tikz@fmgz@tmp@node-m-\the\tikz@fmgz@n);
							\tikz@fmgz@typeout{wiggles final \the\tikz@fmgz@n}%
					 		\draw[shorten >=-0.2pt,shorten <=-0.2pt,propagator] plot[smooth,variable=\t,samples at={0,...,\the\tikz@fmgz@n}]
					  		(\tikz@fmgz@tmp@node-m-\t);
					 		}
						}
				},
				curly wiggles/.style={decoration={%
						markings,
						mark=at position 0 with {\global\tikz@fmgz@n\the\numexpr\pgfkeysvalueof{/tikz/Feynmangraphz/gluon n ini}-1\relax
							\pgfmathtruncatemacro{\tikz@fmgz@k}{round((12*\pgfdecoratedpathlength+1*\pgfkeysvalueof{/tikz/Feynmangraphz/segment length})/(12*\pgfkeysvalueof{/tikz/Feynmangraphz/segment length}))}%	
							\pgfmathtruncatemacro{\tikz@fmgz@nn}{12*\tikz@fmgz@k-2*(\pgfkeysvalueof{/tikz/Feynmangraphz/gluon n ini})+6}% number of steps
							\pgfmathsetmacro{\tikz@fmgz@step}{1/\tikz@fmgz@nn}%
							\xdef\tikz@fmgz@step{\tikz@fmgz@step}%
							\xdef\tikz@fmgz@correctednexpected{\the\numexpr\tikz@fmgz@nn-10}%
							\tikz@fmgz@typeout{path length=\pgfdecoratedpathlength, k=\tikz@fmgz@k, step length=\tikz@fmgz@step, expected step number=\tikz@fmgz@nn}},
						mark=between positions {3*\tikz@fmgz@step} and {1-3*\tikz@fmgz@step} step \tikz@fmgz@step with {
							\global\advance\tikz@fmgz@n1\relax
							\tikz@fmgz@typeout{curly wiggles step=\the\tikz@fmgz@n}%
							\path ({\pgfkeysvalueof{/tikz/Feynmangraphz/amplitude}*sin(\tikz@fmgz@n*30)},
								{\pgfkeysvalueof{/tikz/Feynmangraphz/amplitude}*(cos(\tikz@fmgz@n*30)-0.5)})
								coordinate(\tikz@fmgz@tmp@node-m-\the\tikz@fmgz@n);
							},
						% mark=at position 0.5 with {%\global\advance\tikz@fmgz@n1\relax
						% 	%\path (0,0) coordinate(\tikz@fmgz@tmp@node-m-\the\tikz@fmgz@n);
						% 	%\path (0,{\pgfkeysvalueof{/tikz/Feynmangraphz/amplitude}*sqrt(abs(sin(\tikz@fmgz@n*60)))*sign(sin(\tikz@fmgz@n*60))})
						% 	%	coordinate(\tikz@fmgz@tmp@node-m-\the\tikz@fmgz@n);
						% 	\tikz@fmgz@typeout{curly wiggles final n=\the\tikz@fmgz@n}%
					 	% 	\draw[shorten >=-0.2pt,shorten <=-0.2pt,propagator] plot[smooth,variable=\t,samples at={\pgfkeysvalueof{/tikz/Feynmangraphz/gluon n ini},...,\the\tikz@fmgz@n}]
					 	% 	 (\tikz@fmgz@tmp@node-m-\t);
					 	% 	},
						},	
				},
				draw curly wiggles/.style={
					decoration={%
						show path construction,
						curveto code={%
							\path  (\tikzinputsegmentfirst) coordinate (\tikz@fmgz@tmp@node-m-\the\numexpr\pgfkeysvalueof{/tikz/Feynmangraphz/gluon n ini}\relax);
							\ifnum\tikz@fmgz@n>\tikz@fmgz@correctednexpected\relax
								\path (\tikzinputsegmentlast) coordinate 	(\tikz@fmgz@tmp@node-m-\the\numexpr\tikz@fmgz@n\relax);
							 	\draw[%shorten >=-0.2pt,shorten <=-0.2pt,
							 	-,propagator] plot[smooth,variable=\t,			
							 	samples at={\the\numexpr\pgfkeysvalueof{/tikz/Feynmangraphz/gluon n ini}\relax,...,\the\numexpr\tikz@fmgz@n\relax}]
					 		  	(\tikz@fmgz@tmp@node-m-\t);
							\else
								\typeout{\the\tikz@fmgz@n<=\tikz@fmgz@correctednexpected}%
							\fi	
							},
						lineto code={%
							\path  (\tikzinputsegmentfirst) coordinate (\tikz@fmgz@tmp@node-m-\the\numexpr\pgfkeysvalueof{/tikz/Feynmangraphz/gluon n ini}\relax)
							(\tikzinputsegmentlast) coordinate (\tikz@fmgz@tmp@node-m-\the\numexpr\tikz@fmgz@n\relax);
						 	\draw[shorten >=-0.2pt,shorten <=-0.2pt,-,propagator] plot[smooth,variable=\t,			
							 samples at={\the\numexpr\pgfkeysvalueof{/tikz/Feynmangraphz/gluon n ini}\relax,...,\the\numexpr\tikz@fmgz@n\relax}]
							   (\tikz@fmgz@tmp@node-m-\t);
							},
						moveto code={%
							\typeout{moveto\tikzinputsegmentlast}
							\path (\tikzinputsegmentlast) coordinate (\tikz@fmgz@tmp@node-m-\the\numexpr\tikz@fmgz@n\relax);
							\pgfpathmoveto{\pgfpointanchor{\tikz@fmgz@tmp@node-m-\the\numexpr\tikz@fmgz@n}{center}}%
							\tikzset{color=blue}
							},			
						closepath code=={%
							\typeout{close path}
							},			
							%}
						% markings,
						% mark=at position 0 with {\tikz@fmgz@typeout{curly wiggles final n=\the\tikz@fmgz@n}%
						% 	\draw[shorten >=-0.2pt,shorten <=-0.2pt,-,propagator] plot[smooth,variable=\t,samples at={\pgfkeysvalueof{/tikz/Feynmangraphz/gluon n ini},...,\the\tikz@fmgz@n}]
					 	% 	(\tikz@fmgz@tmp@node-m-\t);}
					 	},
					decorate	
				},
				truly dotted/.style={% truly dotted lines (from https://tex.stackexchange.com/a/52856)
    				dash pattern=on 0.01*\pgflinewidth off ##1*\pgflinewidth,
					line cap=round,
    				shorten >/.evaluated=0.5*\pgflinewidth,
					shorten </.evaluated=0.5*\pgflinewidth
					},
    			truly dotted/.default=6,
				complete dashes/.code args={on ##1 off ##2}{% from https://tex.stackexchange.com/a/133357
					\tikz@addoption{%
						\pgfgetpath\tikz@fmgz@currentpath%
						\pgfprocessround{\tikz@fmgz@currentpath}{\tikz@fmgz@currentpath}%
						\pgf@decorate@parsesoftpath{\tikz@fmgz@currentpath}{\tikz@fmgz@currentpath}%
						\let\tikz@fmgz@len\pgf@decorate@totalpathlength
						\pgfmathtruncatemacro{\tikz@fmgz@nn}{round((\tikz@fmgz@len-##1)/(##1+##2))}%
						\pgfmathsetmacro{\tikz@fmgz@r}{##1/##2}%
						\pgfmathsetmacro{\tikz@fmgz@off}{\tikz@fmgz@len/((\tikz@fmgz@nn+1)*\tikz@fmgz@r+\tikz@fmgz@nn)}%
						\pgfmathsetmacro{\tikz@fmgz@on}{\tikz@fmgz@r*\tikz@fmgz@off}%
			%             \pgfmathparse{\pgf@decorate@totalpathlength-#1}\let\tikz@fmgz@rest=\pgfmathresult%
			%             \pgfmathparse{#1+#2}\let\tikz@fmgz@onoff=\pgfmathresult%
			%             \pgfmathparse{max(floor(\tikz@fmgz@rest/\tikz@fmgz@onoff), 1)}\let\tikz@fmgz@nfullonoff=\pgfmathresult%
			%             \pgfmathparse{max((\tikz@fmgz@rest-\tikz@fmgz@onoff*\tikz@fmgz@nfullonoff)/\tikz@fmgz@nfullonoff+#2, #2)}\let\tikz@fmgz@offexpand=\pgfmathresult%
						\pgfsetdash{{\tikz@fmgz@on}{\tikz@fmgz@off}}{0pt}%
						}%
				},
				complete dashes/.default={on \pgfkeysvalueof{/tikz/Feynmangraphz/dash length} off \pgfkeysvalueof{/tikz/Feynmangraphz/gap length}},
				complete dash cycle/.code args={on ##1 off ##2}{% ends with gap, for loops etc.
					\tikz@addoption{%
						\pgfgetpath\tikz@fmgz@currentpath%
						\pgfprocessround{\tikz@fmgz@currentpath}{\tikz@fmgz@currentpath}%
						\pgf@decorate@parsesoftpath{\tikz@fmgz@currentpath}{\tikz@fmgz@currentpath}%
						\let\tikz@fmgz@len\pgf@decorate@totalpathlength
						\pgfmathtruncatemacro{\tikz@fmgz@nn}{round(\tikz@fmgz@len/(##1+##2))}%
						\pgfmathsetmacro{\tikz@fmgz@r}{##1/##2}%
						\pgfmathsetmacro{\tikz@fmgz@off}{\tikz@fmgz@len/(\tikz@fmgz@nn*\tikz@fmgz@r+\tikz@fmgz@nn)}%
						\pgfmathsetmacro{\tikz@fmgz@on}{\tikz@fmgz@off*\tikz@fmgz@r}%
						\pgfsetdash{{\tikz@fmgz@on}{\tikz@fmgz@off}}{0pt}%
						}%
				},
				complete dash cycle/.default={on \pgfkeysvalueof{/tikz/Feynmangraphz/dash length} off \pgfkeysvalueof{/tikz/Feynmangraphz/gap length}},
				cross/.style={minimum size=1em,path picture={
					\draw[propagator] (path picture bounding box.north west) -- (path picture bounding box.south east)
						(path picture bounding box.north east) -- (path picture bounding box.south west);
				}},
				counter/.style={circle,draw,propagator,minimum size=1em,alias=tmp@counter,append after command={
					 (tmp@counter.45) edge[propagator,line cap=butt,shorten >=0.5\pgflinewidth,shorten <=0.5\pgflinewidth] (tmp@counter.225)
					 (tmp@counter.135) edge[propagator,line cap=butt,shorten >=0.5\pgflinewidth,shorten <=0.5\pgflinewidth] (tmp@counter.315)
				}},
				curved arrow/.style={%
					postaction={Feynmangraphz/curved arrow/set values={##1},
						decoration={
							markings,
							mark=at position 0 with {%
						 		\pgfmathsetmacro{\tikz@fmgz@t}{1.1*\pgfkeysvalueof{/tikz/Feynmangraphz/curved arrow/length}/(\pgfdecoratedpathlength)}
								\xdef\tikz@fmgz@t{\tikz@fmgz@t}},
							mark=at position {\pgfkeysvalueof{/tikz/Feynmangraphz/curved arrow/position}-0.5*\tikz@fmgz@t} with {\coordinate(\tikz@fmgz@tmp@node-0);},
							mark=at position {\pgfkeysvalueof{/tikz/Feynmangraphz/curved arrow/position}-\tikz@fmgz@t/6} with {\coordinate(\tikz@fmgz@tmp@node-1);},
							mark=at position {\pgfkeysvalueof{/tikz/Feynmangraphz/curved arrow/position}+\tikz@fmgz@t/6} with {\coordinate(\tikz@fmgz@tmp@node-2);},
							mark=at position {\pgfkeysvalueof{/tikz/Feynmangraphz/curved arrow/position}+0.5*\tikz@fmgz@t} with {\coordinate(\tikz@fmgz@tmp@node-3);
								\draw[/tikz/Feynmangraphz/curved arrow/arrow keys] %[-\pgfkeysvalueof{/tikz/Feynmangraphz/curved arrow/tip}]
						 			let \p0=(\tikz@fmgz@tmp@node-0),\p1=(\tikz@fmgz@tmp@node-3),\p2=(\tikz@fmgz@tmp@node-1),\p3=(\tikz@fmgz@tmp@node-2),
						  			\n1={(-5*\x0 + 2*\x1 + 18*\x2 - 9*\x3)/6},
						 			\n2={(2*\x0 - 5*\x1 - 9*\x2 + 18*\x3)/6},
						 			\n3={(-5*\y0 + 2*\y1 + 18*\y2 - 9*\y3)/6},
									\n4={(2*\y0 - 5*\y1 - 9*\y2 + 18*\y3)/6}
						 			in
						 			(\tikz@fmgz@tmp@node-0) .. controls (\n1,\n3) and (\n2,\n4) .. (\tikz@fmgz@tmp@node-3);
								},
							},
						decorate}
		 		},
		 		parallel curve/.style={%
			 		postaction={Feynmangraphz/parallel curve/set values={##1},
						decorate,
						decoration={%
						markings, 
						mark=at position 0 with {% 
							\pgfmathsetmacro{\tikz@fmgz@t}{\pgfkeysvalueof{/tikz/Feynmangraphz/parallel	curve/length}/(\pgfdecoratedpathlength)} 
							\xdef\tikz@fmgz@t{\tikz@fmgz@t}}, 
						mark=at position {\pgfkeysvalueof{/tikz/Feynmangraphz/parallel curve/position}-0.5*\tikz@fmgz@t} with 
							{\path (0,\pgfkeysvalueof{/tikz/Feynmangraphz/parallel curve/distance}) coordinate(\tikz@fmgz@tmp@node-0);}, 
						mark=at position {\pgfkeysvalueof{/tikz/Feynmangraphz/parallel curve/position}-\tikz@fmgz@t/6} with
							{\path (0,\pgfkeysvalueof{/tikz/Feynmangraphz/parallel curve/distance}) coordinate(\tikz@fmgz@tmp@node-1);}, 
						mark=at position {\pgfkeysvalueof{/tikz/Feynmangraphz/parallel curve/position}+\tikz@fmgz@t/6} with 
							{\path (0,\pgfkeysvalueof{/tikz/Feynmangraphz/parallel curve/distance}) coordinate(\tikz@fmgz@tmp@node-2);}, 
						mark=at position {\pgfkeysvalueof{/tikz/Feynmangraphz/parallel curve/position}+0.5*\tikz@fmgz@t} with 
							{\path (0,\pgfkeysvalueof{/tikz/Feynmangraphz/parallel curve/distance}) coordinate(\tikz@fmgz@tmp@node-3); 
					 		 \path (\tikz@fmgz@tmp@node-0) edge[Bezier through={(\tikz@fmgz@tmp@node-1)}{(\tikz@fmgz@tmp@node-2)},
								Feynmangraphz/parallel curve/edge options] (\tikz@fmgz@tmp@node-3);}
						}
					},
		 		},
		 		Bezier through/.style 2 args={to path={
			 		let \p0=(\tikztostart),\p1=(\tikztotarget),\p2=##1,\p3=##2,
				  		\n1={(-5*\x0 + 2*\x1 + 18*\x2 - 9*\x3)/6},
				 		\n2={(2*\x0 - 5*\x1 - 9*\x2 + 18*\x3)/6},
				 		\n3={(-5*\y0 + 2*\y1 + 18*\y2 - 9*\y3)/6},
						\n4={(2*\y0 - 5*\y1 - 9*\y2 + 18*\y3)/6}
				 		in
				 		(\tikztostart) .. controls (\n1,\n3) and (\n2,\n4) .. (\tikztotarget) \tikztonodes
		 		}},
		 		Charged/.tip={Stealth[length=\pgfkeysvalueof{/tikz/Feynmangraphz/charged arrow length},
			 		inset=\pgfkeysvalueof{/tikz/Feynmangraphz/charged arrow inset},
					bend]},
		 		FFlow/.tip={Stealth[length=\pgfkeysvalueof{/tikz/Feynmangraphz/fflow arrow length},
			 		inset=\pgfkeysvalueof{/tikz/Feynmangraphz/fflow arrow inset},
					bend]},
		 		Momentum/.tip={Latex[bend]},
		 		momentum arrow/.style={thin,solid,>=Momentum,->},
		 		momentum/.style={parallel curve={##1,edge options/.append style={momentum arrow,momentum annotations,every edge quotes/.append style={inner sep=3pt},##1}}},
		 		momentum annotations/.style={font=\small},
		 		momentum'/.style={parallel curve={distance/.evaluated={-1*(\pgfkeysvalueof{/tikz/Feynmangraphz/parallel curve/distance})*1pt/1cm},
			 		##1,
			 		edge options/.append style={momentum arrow,font=\small,every edge quotes/.append style={inner sep=3pt},swap,##1}}},
		 		%momentum/.search also={/tikz/Feynmangraphz/parallel curve/},
		 		propagator/.style={thick,line cap=round},
		 		propagators/.style={/tikz/propagator/.append style={##1}},
		 		fermion/.style={propagator,solid},
		 		scalar/.style={propagator,complete dashes},
		 		charged scalar/.style={scalar,curved arrow={set values={##1}}},
		 		charged antiscalar/.style={scalar,curved arrow={set values={arrow={<-},##1}}},
		 		%clear to path/.code={\tikzset{to path={}}},
		 		charged fermion/.style={fermion,curved arrow={set values={##1}}},
		 		charged antifermion/.style={fermion,curved arrow={set values={arrow={<-},##1}}},
				ghost/.style={propagator,truly dotted,curved arrow={set values={##1}}},
		 		fermion flow/.style={gray,-FFlow},
		 		scalar/.style={propagator,complete dashes},
		 		% charged scalar/.style={scalar,curved arrow=##1},
		 		% charged scalar/.default=0.5,
		 		arc starting at/.style={to path={let \p1=($(\tikztotarget)-(\tikztostart)$),
					\n1={abs((\x1*\x1+\y1*\y1)/(\x1*cos(##1)+\y1*sin(##1))/2)},
					\n2={scalar(Mod(540-atan2(-2*\x1*\y1*cos(##1)+(\x1*\x1-\y1*\y1)*sin(##1),(\x1*\x1-\y1*\y1)*cos(##1)-2*\x1*\y1*sin(##1)),360))},
					\n3={ifthenelse(\n2>##1,\n2,\n2+360)}
			 		in %\pgfextra{\tikz@fmgz@typeout{start=\tikztostart, target=\tikztotarget, arc[start angle=#1,end angle=\n3,radius=\n1]}}
			 		(\tikztostart) arc[start angle=##1,end angle=\n3,radius=\n1] \tikztonodes
		 		}},
		 		arc' starting at/.style={to path={let \p1=($(\tikztotarget)-(\tikztostart)$),
					\n1={abs((\x1*\x1+\y1*\y1)/(\x1*cos(##1)+\y1*sin(##1))/2)},
					\n2={scalar(Mod(540-atan2(-2*\x1*\y1*cos(##1)+(\x1*\x1-\y1*\y1)*sin(##1),(\x1*\x1-\y1*\y1)*cos(##1)-2*\x1*\y1*sin(##1)),360))},
					\n3={ifthenelse(\n2<##1,\n2,\n2-360)}
			 		in %\pgfextra{\tikz@fmgz@typeout{start=\tikztostart, target=\tikztotarget, arc[start angle=##1,end angle=\n3,radius=\n1]}}
			 		(\tikztostart) arc[start angle=##1,end angle=\n3,radius=\n1] \tikztonodes
				}}
			},
		},
		debug/.is if=@tikz@fmgz@debug,
		every diagram/.style={every edge quotes/.append style={inner sep=.6em}},
		every Feynman node/.style={},
		step count/.code={\advance\c@fmgz@counta by1\relax},
		segment length/.initial=5pt,% for wiggles and gluons
		amplitude/.initial=2pt,% for wiggles and gluons
		aspect/.initial=0,% for wiggles and gluons
        dash length/.initial=3pt,% for complete dashes
        gap length/.initial=4pt,% for complete dashes
		first up/.is if=@tikz@fmgz@firstup,
		layout layer/.code={% 
			\advance\c@fmgz@countb by1\relax
			\def\tikz@fmgz@tmpl##1/##2\pgf@nil{%
			 	\def\tikz@fmgz@tmpnodenameandoptions{##1}%
			 	\def\tikz@fmgz@tmpnodecontents{##2}%
			}%
			\def\tikz@fmgz@tmpm[##1]##2\pgf@nil{%
			 	\def\tikz@fmgz@tmpnodename{##2}%
			 	\def\tikz@fmgz@tmpnodeoptions{##1}%
			}%
			\c@fmgz@counta0\relax
			\tikzset{Feynmangraphz/step count/.list={#1}}%
			\tikz@fmgz@typeout{Found list with \the\c@fmgz@counta\space entries.}%
			\ifnum\c@fmgz@countb=1\relax
				\edef\tikz@fmgz@tmpx{0}%
			\else
				\pgfmathsetmacro{\tikz@fmgz@tmpx}{\tikz@fmgz@tmpx+\pgfkeysvalueof{/tikz/Feynmangraphz/hsep}/sqrt(1+abs(\the\c@fmgz@counta-\tikz@fmgz@tmpn))}%
			\fi
			\edef\tikz@fmgz@tmpn{\the\c@fmgz@counta}%
			\def\pgffor@scan{\pgfutil@ifnextchar({\pgffor@scanround}{\pgfutil@ifnextchar[{\pgffor@scansquare}{\pgffor@scanone}}}%
			\def\pgffor@scansquare[##1]##2,{\def\pgffor@value{[##1]##2}\pgffor@scanned}%
			\foreach \tikz@fmgz@tmpa [count=\tikz@fmgz@tmpb] in {#1}
				{\tikz@fmgz@typeout{set up \tikz@fmgz@tmpa, item number=\tikz@fmgz@tmpb}%
				 %\expanded{\noexpand\pgfutil@in@{/}{\tikz@fmgz@tmpa}}% test if, is contained in \tikz@fmgz@tmpa
				 \expandafter\pgfutil@in@\expandafter/\expandafter{\tikz@fmgz@tmpa}\relax
				 \ifpgfutil@in@
					\expandafter\tikz@fmgz@tmpl\tikz@fmgz@tmpa\pgf@nil\relax%
				 \else
					\let\tikz@fmgz@tmpnodenameandoptions\tikz@fmgz@tmpa
					\let\tikz@fmgz@tmpnodecontents\empty
				 \fi
				 \expandafter\pgfutil@in@\expandafter[\expandafter{\tikz@fmgz@tmpnodenameandoptions}\relax
				 \ifpgfutil@in@
					\expandafter\tikz@fmgz@tmpm\tikz@fmgz@tmpnodenameandoptions\pgf@nil\relax%
					%\tikzset{Feynmangraphz/local node options/.style/.expand once=\tikz@fmgz@tmpnodeoptions}%
				 \else
					 \let\tikz@fmgz@tmpnodename\tikz@fmgz@tmpnodenameandoptions%
					 \let\tikz@fmgz@tmpnodeoptions\empty%
					 %\tikzset{Feynmangraphz/local node options/.style={}}%
				 \fi
				 \expanded{\noexpand\pgfutil@in@{\pgfkeysvalueof{/tikz/Feynmangraphz/vertex prefix}}{\tikz@fmgz@tmpnodename}}% test if the vertex prefix is contained in \tikz@fmgz@tmpnodename
				 \ifpgfutil@in@
					\path (\tikz@fmgz@tmpx,
						{(0.5+0.5*\c@fmgz@counta-\tikz@fmgz@tmpb)*\pgfkeysvalueof{/tikz/Feynmangraphz/vsep}})
						node [/tikz/Feynmangraphz/every Feynman node,%Feynmangraphz/local node options,
							style/.expanded=\tikz@fmgz@tmpnodeoptions,name=\tikz@fmgz@tmpnodename,vertex]{\tikz@fmgz@tmpnodecontents};
				 \else
				 	%\pgfmathsetmacro{\mytmp}{(0.5+0.5*\c@fmgz@counta-\tikz@fmgz@tmpb)*\pgfkeysvalueof{/tikz/Feynmangraphz/vsep}}%
				 	%\tikz@fmgz@typeout{(\tikz@fmgz@tmpx,\mytmp), a=\the\c@fmgz@counta, b=\tikz@fmgz@tmpb}%
					\path (\tikz@fmgz@tmpx,
						{(0.5+0.5*\c@fmgz@counta-\tikz@fmgz@tmpb)*\pgfkeysvalueof{/tikz/Feynmangraphz/vsep}})
						node [/tikz/Feynmangraphz/every Feynman node,%Feynmangraphz/local node options,
							style/.expanded=\tikz@fmgz@tmpnodeoptions,
							name=\tikz@fmgz@tmpnodename] {\tikz@fmgz@tmpnodecontents};
				 \fi
				}
		},
		simple layout/.style={%
			/tikz/execute at begin scope={%
				\c@fmgz@counta0\relax
				\tikzset{Feynmangraphz/step count/.list={#1}}%
				\pgfmathsetmacro{\tikz@fmgz@tmpx}{0}%
				\c@fmgz@countb0\relax
				\tikzset{Feynmangraphz/layout layer/.list={#1}}%
			}
		},
		vertex prefix/.initial=v,
		in label/.initial=i,
		out label/.initial=f,
		vsep/.initial=2,
		hsep/.initial=2,
		curved arrow/.cd,
			length/.initial=10pt,
	 		tip/.initial=Charged,
			position/.initial=0.5,
			set values/.code={\tikzset{Feynmangraphz/curved arrow/.cd,#1}},
			arrow keys/.style={>=\pgfkeysvalueof{/tikz/Feynmangraphz/curved arrow/tip},->},
			arrow/.style={/tikz/Feynmangraphz/curved arrow/arrow keys/.append style={#1}},
		/tikz/Feynmangraphz/.cd,
		gluon n ini/.initial={-4},
		gluon pre/.initial=4pt,
		gluon post/.initial=4pt,
	 	charged arrow tip/.initial=Charged,
	 	charged arrow length/.initial=10pt,
	 	charged arrow inset/.initial=2pt,
	 	fflow arrow tip/.initial=FFlow,
	 	fflow arrow length/.initial=5pt,
	 	fflow arrow inset/.initial=1pt,
		temporary node prefix/.store in=\tikz@fmgz@tmp@node,
		temporary node prefix=fmgz@tmp,
		%parallel curve/.search also={/tikz/Feynmangraphz/parallel curve/edge options},
		parallel curve/.unknown/.code={%
    		\let\searchname=\pgfkeyscurrentname%
			\tikz@fmgz@typeout{unknown key in parallel curve: \searchname}
		},
		parallel curve/.cd,
			ignore/.code={},
			length/.initial=20pt,% <- this is not the length of the emerging curve
			distance/.initial=6pt,
			position/.initial=0.5,
			edge node/.forward to=/tikz/Feynmangraphz/parallel curve/ignore,
			edge options/.style={length/.forward to=/tikz/Feynmangraphz/parallel curve/length,
				distance/.forward to=/tikz/Feynmangraphz/parallel curve/distance,
				position/.forward to=/tikz/Feynmangraphz/parallel curve/position,
			},
			edge options+/.style={/tikz/Feynmangraphz/parallel curve/.append style={#1}},
			%edge options/.search also=/tikz/Feynmangraphz/parallel curve/,
			style/.style={},
% 			has label/.is if=tikz@fmgz@paralell@curve@has@label,
% 			label/.code={\tikz@fmgz@typeout{label!}\tikz@fmgz@paralell@curve@has@labeltrue
% 				\def\tikz@fmgz@paralell@curve@label{#1}},
% 			parallel label/.style={midway,font=\tiny,/tikz/Feynmangraphz/parallel curve/distant dependent},
			distant dependent/.code={\ifdim\pgfkeysvalueof{/tikz/Feynmangraphz/parallel curve/distance}>0pt
				\tikzset{allow upside down,below}
			\else
				\tikzset{allow upside down,above}
			\fi},
			has label=false,
			set values/.code={\tikzset{Feynmangraphz/parallel curve/.cd,#1}}
}
\newcount\tikz@fmgz@n
\tikzset{even odd clip/.code={\pgfseteorule},% from https://tex.stackexchange.com/a/538932
	}

% \def\pgfpoint@on@fmgz@gluon#1#2#3{%
% 	\pgf@x=#1\pgfdecorationsegmentamplitude%
% 	\pgf@x=\pgfdecorationsegmentaspect\pgf@x%
% 	\pgf@y=#2\pgfdecorationsegmentamplitude%
% 	\pgf@xa=0.083333333333\pgfdecorationsegmentlength%
% 	\advance\pgf@x by#3\pgf@xa%
% 	%\advance\pgf@x by-\generaloffset pt%
%   }	

\pgfdeclaredecoration{fmgz photon}{initial}{%
	\state{initial}[persistent precomputation={%
			\tikzset{/pgf/fpu/install only={reciprocal}}%
			\global\tikz@fmgz@n0\relax
			\pgfmathtruncatemacro{\tikz@fmgz@nn}{round((\pgfdecoratedpathlength)/\pgfdecorationsegmentlength)}% number of steps
			\ifodd\tikz@fmgz@nn
				\edef\tikz@fmgz@nn{\the\numexpr\tikz@fmgz@nn+1}%
			\fi
			\pgfmathsetmacro{\tikz@fmgz@step}{(\pgfdecoratedpathlength-0.01pt)/\tikz@fmgz@nn}% adjusted segment length
			\pgfdecorationsegmentlength=\tikz@fmgz@step pt%
			\edef\tikz@fmgz@correctednexpected{\the\numexpr\tikz@fmgz@nn-2}%
			\tikz@fmgz@typeout{path length=\pgfdecoratedpathlength, step length=\tikz@fmgz@step, expected step number=\tikz@fmgz@nn}%	
			\pgfmathsetmacro{\tikz@fmgz@r}{0.45*min(\pgfdecorationsegmentlength,\pgfdecorationsegmentamplitude)}%		
			\pgfset{inner sep=0pt,minimum width=\tikz@fmgz@r pt}%
			\tikz@fmgz@typeout{r=\tikz@fmgz@r pt, segment length=\the\pgfdecorationsegmentlength, amplitude=\the\pgfdecorationsegmentamplitude}%
	},width=0pt,next state=wiggle]{%
		% \pgfpathmoveto{\pgfpointdecoratedpathfirst}%
		% \pgfpathmoveto{\pgfpoint{\pgfdecorationsegmentlength}{0pt}}%
		% \pgfinterruptpath
		%   \pgfnode{circle}{center}{}{\tikz@fmgz@tmp@node-m-\the\tikz@fmgz@n}{\pgfusepath{stroke}}%
		% \endpgfinterruptpath
	}%
	\state{wiggle}[width=\tikz@fmgz@step pt,switch if less than={\pgfdecorationsegmentlength} to final]{%
		%\pgfpointanchor{\tikz@fmgz@tmp@node-m-\the\tikz@fmgz@n}{center}}%
		\global\advance\tikz@fmgz@n1\relax
		%\ifnum\tikz@fmgz@n=1\relax
		\ifodd\tikz@fmgz@n
			\pgfpathcurveto%
				{\pgfpoint{0.25\pgfdecorationsegmentlength}{0.5\pgfdecorationsegmentamplitude}}%
				{\pgfpoint{0.75\pgfdecorationsegmentlength}{0.5\pgfdecorationsegmentamplitude}}%
				{\pgfpoint{\pgfdecorationsegmentlength}{0pt}}%	
		\else
			\pgfpathcurveto%
				{\pgfpoint{0.25\pgfdecorationsegmentlength}{-0.5\pgfdecorationsegmentamplitude}}%
				{\pgfpoint{0.75\pgfdecorationsegmentlength}{-0.5\pgfdecorationsegmentamplitude}}%
				{\pgfpoint{\pgfdecorationsegmentlength}{0pt}}%	
		\fi
    }%
	\state{final}[width=1sp]{%
		%\pgfpathmoveto{\pgfpointshapeborder{\tikz@fmgz@tmp@node-m-\the\numexpr\tikz@fmgz@n\relax}{\pgfpointadd{\pgfpointanchor{\tikz@fmgz@tmp@node-m-\the\numexpr\tikz@fmgz@n\relax}{center}}{\pgfpoint{-0.4\pgfdecorationsegmentlength}{\pgfdecorationsegmentamplitude}}}}%
		%\pgfpathlineto{\pgfpointdecoratedpathlast}%
		%\pgfpathcurveto%
		%		{\pgfpoint{-0.5\pgfdecorationsegmentlength}{-0.5\pgfdecorationsegmentamplitude}}%
		%		{\pgfpoint{-0.1\pgfdecorationsegmentlength}{0\pgfdecorationsegmentamplitude}}%
		%		{\pgfpointdecoratedpathlast}%
	}%
}	

\pgfdeclaredecoration{fmgz photon'}{initial}{%
	\state{initial}[persistent precomputation={%
			\tikzset{/pgf/fpu/install only={reciprocal}}%
			\global\tikz@fmgz@n0\relax
			\pgfmathtruncatemacro{\tikz@fmgz@nn}{round((\pgfdecoratedpathlength)/\pgfdecorationsegmentlength)}% number of steps
			\ifodd\tikz@fmgz@nn
				\edef\tikz@fmgz@nn{\the\numexpr\tikz@fmgz@nn+1}%
			\fi
			\pgfmathsetmacro{\tikz@fmgz@step}{(\pgfdecoratedpathlength-0.01pt)/\tikz@fmgz@nn}% adjusted segment length
			\pgfdecorationsegmentlength=\tikz@fmgz@step pt%
			\edef\tikz@fmgz@correctednexpected{\the\numexpr\tikz@fmgz@nn-2}%
			\tikz@fmgz@typeout{path length=\pgfdecoratedpathlength, step length=\tikz@fmgz@step, expected step number=\tikz@fmgz@nn}%	
			\pgfmathsetmacro{\tikz@fmgz@r}{0.45*min(\pgfdecorationsegmentlength,\pgfdecorationsegmentamplitude)}%		
			\pgfset{inner sep=0pt,minimum width=\tikz@fmgz@r pt}%
			\tikz@fmgz@typeout{r=\tikz@fmgz@r pt, segment length=\the\pgfdecorationsegmentlength, amplitude=\the\pgfdecorationsegmentamplitude}%
	},width=0pt,next state=wiggle]{%
	}%
	\state{wiggle}[width=\tikz@fmgz@step pt,switch if less than={\pgfdecorationsegmentlength} to final]{%
		\global\advance\tikz@fmgz@n1\relax
		\unless\ifodd\tikz@fmgz@n
			\pgfpathcurveto%
				{\pgfpoint{0.25\pgfdecorationsegmentlength}{0.5\pgfdecorationsegmentamplitude}}%
				{\pgfpoint{0.75\pgfdecorationsegmentlength}{0.5\pgfdecorationsegmentamplitude}}%
				{\pgfpoint{\pgfdecorationsegmentlength}{0pt}}%	
		\else
			\pgfpathcurveto%
				{\pgfpoint{0.25\pgfdecorationsegmentlength}{-0.5\pgfdecorationsegmentamplitude}}%
				{\pgfpoint{0.75\pgfdecorationsegmentlength}{-0.5\pgfdecorationsegmentamplitude}}%
				{\pgfpoint{\pgfdecorationsegmentlength}{0pt}}%	
		\fi
    }%
	\state{final}[width=1sp]{%
	}%
}	

\pgfdeclaredecoration{fmgz gluon}{initial}{%
	\state{initial}[persistent precomputation={%
			\tikzset{/pgf/fpu/install only={reciprocal}}%
			\global\tikz@fmgz@n0\relax
			\pgfmathtruncatemacro{\tikz@fmgz@nn}{round((\pgfdecoratedpathlength-0.95*\pgfdecorationsegmentlength)/\pgfdecorationsegmentlength)}% number of steps
			\pgfmathsetmacro{\tikz@fmgz@step}{(\pgfdecoratedpathlength-0.95*\pgfdecorationsegmentlength-1pt)/\tikz@fmgz@nn}% adjusted segment length
			\pgfdecorationsegmentlength=\tikz@fmgz@step pt%
			\edef\tikz@fmgz@correctednexpected{\the\numexpr\tikz@fmgz@nn-2}%
			\tikz@fmgz@typeout{path length=\pgfdecoratedpathlength, step length=\tikz@fmgz@step, expected step number=\tikz@fmgz@nn}%	
			\pgfmathsetmacro{\tikz@fmgz@r}{0.45*min(\pgfdecorationsegmentlength,\pgfdecorationsegmentamplitude)}%		
			\pgfset{inner sep=0pt,minimum width=\tikz@fmgz@r pt}%
			\tikz@fmgz@typeout{r=\tikz@fmgz@r pt, segment length=\the\pgfdecorationsegmentlength, amplitude=\the\pgfdecorationsegmentamplitude}%
	},width=0.95\pgfdecorationsegmentlength,next state=coil]{%
		% \pgfpathmoveto{\pgfpointdecoratedpathfirst}%
		% \pgfpathmoveto{\pgfpoint{\pgfdecorationsegmentlength}{0pt}}%
		% \pgfinterruptpath
		%   \pgfnode{circle}{center}{}{\tikz@fmgz@tmp@node-m-\the\tikz@fmgz@n}{\pgfusepath{stroke}}%
		% \endpgfinterruptpath
	}%
	\state{coil}[width=\tikz@fmgz@step pt,switch if less than={\pgfdecorationsegmentlength} to final]{%
		%\pgfpointanchor{\tikz@fmgz@tmp@node-m-\the\tikz@fmgz@n}{center}}%
		\global\advance\tikz@fmgz@n1\relax
		\ifnum\tikz@fmgz@n=1\relax
			\pgfinterruptpath
				\pgfnode{circle}{center}{}{\tikz@fmgz@tmp@node-m-\the\tikz@fmgz@n}{}%
			\endpgfinterruptpath
			\pgfpathmoveto{\pgfpointdecoratedpathfirst}%
			\pgfpathcurveto%
				{\pgfpoint{-0.5\pgfdecorationsegmentlength}{0.5\pgfdecorationsegmentamplitude}}%
				{\pgfpoint{-0.2\pgfdecorationsegmentlength}{0.5\pgfdecorationsegmentamplitude}}%
				{\pgfpointshapeborder{\tikz@fmgz@tmp@node-m-\the\numexpr\tikz@fmgz@n\relax}{\pgfpointadd{\pgfpointanchor{\tikz@fmgz@tmp@node-m-\the\numexpr\tikz@fmgz@n\relax}{center}}{\pgfpoint{-0.4\pgfdecorationsegmentlength}{\pgfdecorationsegmentamplitude}}}}%	
		\else
			\pgfpathmoveto{\pgfpoint{\pgfdecorationsegmentlength}{0pt}}%
			\pgfinterruptpath
			\pgfnode{circle}{center}{}{\tikz@fmgz@tmp@node-m-\the\tikz@fmgz@n}{}%
			\endpgfinterruptpath
			%\pgfpathcircle{\pgfpointshapeborder{\tikz@fmgz@tmp@node-m-\the\numexpr\tikz@fmgz@n-1\relax}{\pgfpointadd{\pgfpointanchor{\tikz@fmgz@tmp@node-m-\the\numexpr\tikz@fmgz@n-1\relax}{center}}{\pgfpoint{0.1\pgfdecorationsegmentlength}{-\pgfdecorationsegmentamplitude}}}}{0.2*\tikz@fmgz@r pt}%
			\pgfpathmoveto{\pgfpointshapeborder{\tikz@fmgz@tmp@node-m-\the\numexpr\tikz@fmgz@n-1\relax}{\pgfpointadd{\pgfpointanchor{\tikz@fmgz@tmp@node-m-\the\numexpr\tikz@fmgz@n-1\relax}{center}}{\pgfpoint{0.1\pgfdecorationsegmentlength}{-\pgfdecorationsegmentamplitude}}}}%
			\pgfpathcurveto%
				{\pgfpoint{-0.8\pgfdecorationsegmentlength}{-\pgfdecorationsegmentamplitude}}%
				{\pgfpoint{-2\pgfdecorationsegmentlength}{-\pgfdecorationsegmentamplitude}}%
				{\pgfpointanchor{\tikz@fmgz@tmp@node-m-\the\numexpr\tikz@fmgz@n-1\relax}{center}}%
			\pgfpathcurveto%
				{\pgfpoint{-0.5\pgfdecorationsegmentlength}{0.5\pgfdecorationsegmentamplitude}}%
				{\pgfpoint{-0.2\pgfdecorationsegmentlength}{0.5\pgfdecorationsegmentamplitude}}%
				{\pgfpointshapeborder{\tikz@fmgz@tmp@node-m-\the\numexpr\tikz@fmgz@n\relax}{\pgfpointadd{\pgfpointanchor{\tikz@fmgz@tmp@node-m-\the\numexpr\tikz@fmgz@n\relax}{center}}{\pgfpoint{-0.4\pgfdecorationsegmentlength}{\pgfdecorationsegmentamplitude}}}}%			
			\tikz@fmgz@typeout{x=\the\pgf@x, y=\the\pgf@y, r=\tikz@fmgz@r pt, segment length=\the\pgfdecorationsegmentlength, amplitude=\the\pgfdecorationsegmentamplitude}%
		\fi
    }%
	\state{final}[width=1sp]{%
		\pgfpathmoveto{\pgfpointshapeborder{\tikz@fmgz@tmp@node-m-\the\numexpr\tikz@fmgz@n\relax}{\pgfpointadd{\pgfpointanchor{\tikz@fmgz@tmp@node-m-\the\numexpr\tikz@fmgz@n\relax}{center}}{\pgfpoint{-0.4\pgfdecorationsegmentlength}{\pgfdecorationsegmentamplitude}}}}%
		%\pgfpathlineto{\pgfpointdecoratedpathlast}%
		\pgfpathcurveto%
				{\pgfpoint{-0.5\pgfdecorationsegmentlength}{-0.5\pgfdecorationsegmentamplitude}}%
				{\pgfpoint{-0.1\pgfdecorationsegmentlength}{0\pgfdecorationsegmentamplitude}}%
				{\pgfpointdecoratedpathlast}%
	}%
}	


% \pgfdeclaredecoration{fmgz gluon}{initial}{%
% 	\state{initial}[width=0pt,next state=mark,persistent precomputation=\tikzset{/pgf/fpu/install only={reciprocal}}] {%
% 		\global\tikz@fmgz@n\the\numexpr\pgfkeysvalueof{/tikz/Feynmangraphz/gluon n ini}\relax
% 		\pgfinterruptpath%
% 		  \pgfscope%
% 			\pgfcoordinate{\tikz@fmgz@tmp@node-m-\the\tikz@fmgz@n}{\pgfpointdecoratedpathfirst}%
% 			\pgfcoordinate{\tikz@fmgz@tmp@node-m-last}{\pgfpointdecoratedpathlast}%
% 		  \endpgfscope%
% 		\endpgfinterruptpath%
% 		\pgfmathtruncatemacro{\tikz@fmgz@k}{round((12*\pgfdecoratedpathlength+1*\pgfkeysvalueof{/tikz/Feynmangraphz/segment length})/(12*\pgfkeysvalueof{/tikz/Feynmangraphz/segment length}))}%	
% 		\pgfmathtruncatemacro{\tikz@fmgz@nn}{12*\tikz@fmgz@k-2*(\pgfkeysvalueof{/tikz/Feynmangraphz/gluon n ini})-2}% number of steps
% 		\pgfmathsetmacro{\tikz@fmgz@step}{\pgfdecoratedpathlength/\tikz@fmgz@nn}%
% 		\xdef\tikz@fmgz@step{\tikz@fmgz@step}%
% 		\xdef\tikz@fmgz@correctednexpected{\the\numexpr\tikz@fmgz@nn-2}%
% 		\tikz@fmgz@typeout{path length=\pgfdecoratedpathlength, k=\tikz@fmgz@k, step length=\tikz@fmgz@step, expected step number=\tikz@fmgz@nn}%
% 	}
% 	\state{mark}[width=\tikz@fmgz@step pt,switch if less than={\tikz@fmgz@step} to final]{%
% 		%\ifnum\tikz@fmgz@step>0\relax
% 		\global\advance\tikz@fmgz@n1\relax
% 		\pgfmathtruncatemacro{\tikz@fmgz@itest}{(\tikz@fmgz@n-\pgfkeysvalueof{/tikz/Feynmangraphz/gluon n ini}>4)&&(\tikz@fmgz@correctednexpected-\tikz@fmgz@n>4)}%
% 		\tikz@fmgz@typeout{step \the\tikz@fmgz@n, test=\tikz@fmgz@itest}%
% 		\pgfinterruptpath%
% 		  \pgfscope%
% 		    \ifnum\tikz@fmgz@itest=1\relax
% 				\pgfcoordinate{\tikz@fmgz@tmp@node-m-\the\tikz@fmgz@n}{\pgfpoint{\pgfkeysvalueof{/tikz/Feynmangraphz/amplitude}*(1-0.2*cos(\tikz@fmgz@n*30))*sin(\tikz@fmgz@n*30)}{\pgfkeysvalueof{/tikz/Feynmangraphz/amplitude}*(cos(\tikz@fmgz@n*30)-0)}}%
% 			\else
% 				\pgfcoordinate{\tikz@fmgz@tmp@node-m-\the\tikz@fmgz@n}{\pgfpoint{0pt}{0pt}}%
% 			\fi	
% 		  \endpgfscope%
% 		\endpgfinterruptpath%
% 		%\fi    
% 		}
% 	\state{final}{%
% 		\global\advance\tikz@fmgz@n1\relax
% 		\pgfcoordinate{\tikz@fmgz@tmp@node-m-\the\tikz@fmgz@n}{\pgfpointanchor{\tikz@fmgz@tmp@node-m-last}{center}}%
% 		\draw[%shorten >=-0.2pt,shorten <=-0.2pt,
% 			-,propagator] plot[smooth,variable=\t,			
% 			samples at={\the\numexpr\pgfkeysvalueof{/tikz/Feynmangraphz/gluon n ini}\relax,%
% 			\the\numexpr\pgfkeysvalueof{/tikz/Feynmangraphz/gluon n ini}+2\relax,
% 			\the\numexpr\pgfkeysvalueof{/tikz/Feynmangraphz/gluon n ini}+3\relax,
% 			...,
% 			%\the\numexpr\tikz@fmgz@n-2\relax,
% 			\the\numexpr\tikz@fmgz@n\relax}] (\tikz@fmgz@tmp@node-m-\t);
% 	   %\draw[red] (\tikz@fmgz@tmp@node-m-\the\numexpr\pgfkeysvalueof{/tikz/Feynmangraphz/gluon n ini})   -- (\tikz@fmgz@tmp@node-m-last); 
% 		%\pgfpathlineto{\pgfpointdecoratedpathlast}
% 	} 
% }	
	
% obsolete or to become obsolete

% the challenge the following approach has is that the nonlinear transformations may lead to very uneven wiggle lengths	
	% 	wiggles/.style={decoration={show path construction,
	% 		moveto code={},
	%     	lineto code={\expanded{\noexpand\tikz@fmgz@pgfdist\tikzinputsegmentfirst;\tikzinputsegmentlast\noexpand\pgf@stop}%
	% 		  \let\tikz@fmgz@len\pgfmathresult
	% 		  \expanded{\noexpand\tikz@fmgz@pgfangle\tikzinputsegmentfirst;\tikzinputsegmentlast\noexpand\pgf@stop}%
	% 		  \let\tikz@fmgz@angle\pgfmathresult
	% 		  \pgfmathtruncatemacro{\tikz@fmgz@nn}{round(\tikz@fmgz@len/\pgfkeysvalueof{/tikz/Feynmangraphz/segment length})}%
	% 		  \pgfmathsetmacro{\tikz@fmgz@segementlength}{\tikz@fmgz@len/\tikz@fmgz@nn/1cm}%
	% 		  \if@tikz@fmgz@firstup
	% 	  		\draw[shift={(\tikzinputsegmentfirst)},rotate=\tikz@fmgz@angle,propagator]
	% 				plot[smooth,variable=\t,domain=0:\tikz@fmgz@nn,samples=12*\tikz@fmgz@nn+1]
	% 				({\tikz@fmgz@segementlength*\t},
	% 				 {\pgfkeysvalueof{/tikz/Feynmangraphz/amplitude}*sqrt(abs(sin(360*\t)))*sign(sin(360*\t))});
	% 			\global\@tikz@fmgz@firstupfalse
	% 		  \else
	% 	  		\draw[shift={(\tikzinputsegmentfirst)},rotate=\tikz@fmgz@angle,propagator]
	% 				plot[smooth,variable=\t,domain=0:\tikz@fmgz@nn,samples=12*\tikz@fmgz@nn+1]
	% 				({\tikz@fmgz@segementlength*\t},
	% 				{-\pgfkeysvalueof{/tikz/Feynmangraphz/amplitude}*sqrt(abs(sin(360*\t)))*sign(sin(360*\t))});
	% 			\global\@tikz@fmgz@firstuptrue
	% 		  \fi
	%     	},
	%     	curveto code={
	% 			%\expanded{\noexpand\tikz@fmgz@pgfdist\tikzinputsegmentfirst;\tikzinputsegmentlast\noexpand\pgf@stop}%
	% 			%\let\tikz@fmgz@len\pgfmathresult
	% 			\pgfmathsetmacro{\tikz@fmgz@len}{\pgfdecoratedpathlength}%
	% 			\tikzset{allow upside down}%
	% 		    \pgfsetcurvilinearbeziercurve
	%     			{\expandafter\tikz@fmgz@pgfpoint\tikzinputsegmentfirst\pgf@stop}%
	%     			{\expandafter\tikz@fmgz@pgfpoint\tikzinputsegmentsupporta\pgf@stop}%
	%     			{\expandafter\tikz@fmgz@pgfpoint\tikzinputsegmentsupportb\pgf@stop}%
	%     			{\expandafter\tikz@fmgz@pgfpoint\tikzinputsegmentlast\pgf@stop}%
	% 		  \pgftransformnonlinear{\pgfpointcurvilinearbezierorthogonal\pgf@x\pgf@y}%
	% 		  \expanded{\noexpand\pgfcurvilineardistancetotime{\tikz@fmgz@len pt}}%
	% 		  \pgfmathsetmacro{\tikz@fmgz@len}{\tikz@fmgz@len/\pgf@x}%
	% 		  %\let\tikz@fmgz@len\pgfdecoratedpathlength
	% 		  \pgfmathtruncatemacro{\tikz@fmgz@nn}{round(\tikz@fmgz@len/\pgfkeysvalueof{/tikz/Feynmangraphz/segment length})}%
	% 		  \pgfmathsetmacro{\tikz@fmgz@segementlength}{\tikz@fmgz@len/\tikz@fmgz@nn/1cm}%
	% 		  \if@tikz@fmgz@firstup
	% 	  		\draw[propagator] plot[smooth,variable=\t,domain=0:\tikz@fmgz@nn,samples=12*\tikz@fmgz@nn+1]
	% 				({\tikz@fmgz@segementlength*\t},
	% 				 {\pgfkeysvalueof{/tikz/Feynmangraphz/amplitude}*sqrt(abs(sin(360*\t)))*sign(sin(360*\t))});
	% 			%\global\@tikz@fmgz@firstupfalse
	% 			%\draw[{Circle}-{Circle},blue] (0,0) -- (\tikz@fmgz@len pt,0);
	% 		  \else
	% 	  		\draw[propagator] plot[smooth,variable=\t,domain=0:\tikz@fmgz@nn,samples=12*\tikz@fmgz@nn+1]
	% 				({\tikz@fmgz@segementlength*\t},
	% 				{-\pgfkeysvalueof{/tikz/Feynmangraphz/amplitude}*sqrt(abs(sin(360*\t)))*sign(sin(360*\t))});
	% 			%\global\@tikz@fmgz@firstuptrue
	% 			%\draw[{Circle}-{Circle},blue] (0,0) -- (\tikz@fmgz@len pt,0);
	% 		  \fi
	%     	},
	%     	closepath code={}
	% 		}
	% 	},
	% 	gluon/.style={decoration={show path construction,
	% 		moveto code={},
	%     	lineto code={\expanded{\noexpand\tikz@fmgz@pgfdist\tikzinputsegmentfirst;\tikzinputsegmentlast\noexpand\pgf@stop}%
	% 		  \let\tikz@fmgz@len\pgfmathresult
	% 		  \expanded{\noexpand\tikz@fmgz@pgfangle\tikzinputsegmentfirst;\tikzinputsegmentlast\noexpand\pgf@stop}%
	% 		  \let\tikz@fmgz@angle\pgfmathresult
	% 		  \pgfmathtruncatemacro{\tikz@fmgz@nn}{round(\tikz@fmgz@len/\pgfkeysvalueof{/tikz/Feynmangraphz/segment length})}%
	% 		  \pgfmathsetmacro{\tikz@fmgz@segementlength}{\tikz@fmgz@len/(\tikz@fmgz@nn+0.5)/1cm}%
	% 		  \if@tikz@fmgz@firstup
	% 	  		\draw[shift={(\tikzinputsegmentfirst)},rotate=\tikz@fmgz@angle,propagator] plot[smooth,variable=\t,domain=0:\tikz@fmgz@nn,samples=12*\tikz@fmgz@nn+1]
	% 				({\tikz@fmgz@segementlength*\t+\tikz@fmgz@segementlength*\pgfkeysvalueof{/tikz/Feynmangraphz/aspect}*(1-cos(360*\t))},
	% 				 {2*\pgfkeysvalueof{/tikz/Feynmangraphz/amplitude}*sin(360*\t)});
	% 		  \else
	% 	  		\draw[shift={(\tikzinputsegmentfirst)},rotate=\tikz@fmgz@angle,propagator] plot[smooth,variable=\t,domain=0:\tikz@fmgz@nn,samples=12*\tikz@fmgz@nn+1]
	% 				({\tikz@fmgz@segementlength*\t+\tikz@fmgz@segementlength*\pgfkeysvalueof{/tikz/Feynmangraphz/aspect}*(1-cos(360*\t))},
	% 				{-2*\pgfkeysvalueof{/tikz/Feynmangraphz/amplitude}*sin(360*\t)});
	% 		  \fi
	%     	},
	%     	curveto code={\expanded{\noexpand\tikz@fmgz@pgfdist\tikzinputsegmentfirst;\tikzinputsegmentlast\noexpand\pgf@stop}%
	% 			\let\tikz@fmgz@len\pgfmathresult
	% 			\path (\tikzinputsegmentfirst) coordinate (tmp);
	% 		  \pgfsetcurvilinearbeziercurve
	%     			{\expandafter\tikz@fmgz@pgfpoint\tikzinputsegmentfirst\pgf@stop}%
	%     			{\expandafter\tikz@fmgz@pgfpoint\tikzinputsegmentsupporta\pgf@stop}%
	%     			{\expandafter\tikz@fmgz@pgfpoint\tikzinputsegmentsupportb\pgf@stop}%
	%     			{\expandafter\tikz@fmgz@pgfpoint\tikzinputsegmentlast\pgf@stop}%
	% 		  \pgftransformnonlinear{\pgfpointcurvilinearbezierorthogonal\pgf@x\pgf@y}%
	% 		  \expanded{\noexpand\pgfcurvilineardistancetotime{\tikz@fmgz@len pt}}%
	% 		  \pgfmathsetmacro{\tikz@fmgz@len}{\tikz@fmgz@len/\pgf@x}%
	% 		  \pgfmathtruncatemacro{\tikz@fmgz@nn}{round(\tikz@fmgz@len/\pgfkeysvalueof{/tikz/Feynmangraphz/segment length})}%
	% 		  \pgfmathsetmacro{\tikz@fmgz@segementlength}{\tikz@fmgz@len/(\tikz@fmgz@nn+0)/1cm}%
	% 		  \if@tikz@fmgz@firstup
	% 	  		\draw[propagator] plot[smooth,variable=\t,domain=0:\tikz@fmgz@nn,samples=12*\tikz@fmgz@nn+1]
	% 				({\tikz@fmgz@segementlength*\t+\tikz@fmgz@segementlength*\pgfkeysvalueof{/tikz/Feynmangraphz/aspect}*(1-cos(360*\t))},
	% 				 {2*\pgfkeysvalueof{/tikz/Feynmangraphz/amplitude}*sin(360*\t)});
	% 		  \else
	% 	  		\draw[propagator] plot[smooth,variable=\t,domain=0:\tikz@fmgz@nn,samples=12*\tikz@fmgz@nn+1]
	% 				({\tikz@fmgz@segementlength*\t+\tikz@fmgz@segementlength*\pgfkeysvalueof{/tikz/Feynmangraphz/aspect}*(1-cos(360*\t))},
	% 				{-2*\pgfkeysvalueof{/tikz/Feynmangraphz/amplitude}*sin(360*\t)});
	% 		  \fi
	%     	},
	%     	closepath code={}
	% 		}
	%     },
	%     complete dashes/.style={decorate,decoration={show path construction,
	%   		moveto code={},
	%       	lineto code={\expanded{\noexpand\tikz@fmgz@pgfdist\tikzinputsegmentfirst;\tikzinputsegmentlast\noexpand\pgf@stop}%
	%   		  \let\tikz@fmgz@len\pgfmathresult
	%   		 %\expanded{\noexpand\tikz@fmgz@pgfangle\tikzinputsegmentfirst;\tikzinputsegmentlast\noexpand\pgf@stop}%
	%   		  %\let\tikz@fmgz@angle\pgfmathresult
	%         \pgfmathtruncatemacro{\tikz@fmgz@nn}{round((\tikz@fmgz@len-\pgfkeysvalueof{/tikz/Feynmangraphz/dash length})/(\pgfkeysvalueof{/tikz/Feynmangraphz/dash length}+\pgfkeysvalueof{/tikz/Feynmangraphz/gap length}))}%
	%         \pgfmathsetmacro{\tikz@fmgz@r}{\pgfkeysvalueof{/tikz/Feynmangraphz/dash length}/\pgfkeysvalueof{/tikz/Feynmangraphz/gap length}}%
	%         \pgfmathsetmacro{\tikz@fmgz@off}{\tikz@fmgz@len/((\tikz@fmgz@nn+1)*\tikz@fmgz@r+\tikz@fmgz@nn)}%
	%         \pgfmathsetmacro{\tikz@fmgz@on}{\tikz@fmgz@r*\tikz@fmgz@off}%
	%   		  %\pgfmathtruncatemacro{\tikz@fmgz@nn}{round(\tikz@fmgz@len/\pgfkeysvalueof{/tikz/Feynmangraphz/segment length})}%
	%   		  %\pgfmathsetmacro{\tikz@fmgz@segementlength}{\tikz@fmgz@len/\tikz@fmgz@nn/1cm}%
	%         \tikz@fmgz@typeout{length=\tikz@fmgz@len, n=\tikz@fmgz@nn, on=\tikz@fmgz@on, off=\tikz@fmgz@off}%
	%   	  	\draw[propagator,dash pattern=on \tikz@fmgz@on pt off \tikz@fmgz@off pt]
	%   				(\tikzinputsegmentfirst) -- (\tikzinputsegmentlast);
	%       	},
	%       	curveto code={\expanded{\noexpand\tikz@fmgz@pgfdist\tikzinputsegmentfirst;\tikzinputsegmentlast\noexpand\pgf@stop}%
	%   			\let\tikz@fmgz@len\pgfmathresult
	%   			\tikzset{allow upside down}%
	%   		  \pgfsetcurvilinearbeziercurve
	%       			{\expandafter\tikz@fmgz@pgfpoint\tikzinputsegmentfirst\pgf@stop}%
	%       			{\expandafter\tikz@fmgz@pgfpoint\tikzinputsegmentsupporta\pgf@stop}%
	%       			{\expandafter\tikz@fmgz@pgfpoint\tikzinputsegmentsupportb\pgf@stop}%
	%       			{\expandafter\tikz@fmgz@pgfpoint\tikzinputsegmentlast\pgf@stop}%
	%   		  \pgftransformnonlinear{\pgfpointcurvilinearbezierorthogonal\pgf@x\pgf@y}%
	%   		  \expanded{\noexpand\pgfcurvilineardistancetotime{\tikz@fmgz@len pt}}%
	%   		  \pgfmathsetmacro{\tikz@fmgz@len}{\tikz@fmgz@len/\pgf@x}%
	%         \pgfmathtruncatemacro{\tikz@fmgz@nn}{round((\tikz@fmgz@len-\pgfkeysvalueof{/tikz/Feynmangraphz/dash length})/(\pgfkeysvalueof{/tikz/Feynmangraphz/dash length}+\pgfkeysvalueof{/tikz/Feynmangraphz/gap length}))}%
	%         \pgfmathsetmacro{\tikz@fmgz@r}{\pgfkeysvalueof{/tikz/Feynmangraphz/dash length}/\pgfkeysvalueof{/tikz/Feynmangraphz/gap length}}%
	%         \pgfmathsetmacro{\tikz@fmgz@off}{\tikz@fmgz@len/((\tikz@fmgz@nn+1)*\tikz@fmgz@r+\tikz@fmgz@nn)}%
	%         \pgfmathsetmacro{\tikz@fmgz@on}{\tikz@fmgz@r*\tikz@fmgz@off}%
	%   		  %\pgfmathtruncatemacro{\tikz@fmgz@nn}{round(\tikz@fmgz@len/\pgfkeysvalueof{/tikz/Feynmangraphz/segment length})}%
	%   		  %\pgfmathsetmacro{\tikz@fmgz@segementlength}{\tikz@fmgz@len/\tikz@fmgz@nn/1cm}%
	%         \tikz@fmgz@typeout{length=\tikz@fmgz@len, n=\tikz@fmgz@nn, on=\tikz@fmgz@on, off=\tikz@fmgz@off}%
	%   	  	\draw[propagator,dash pattern=on \tikz@fmgz@on pt off \tikz@fmgz@off pt]
	%   				(0,0) -- (\tikz@fmgz@len pt,0);
	%         },
	%       	closepath code={}
	%   		}
	%     },
	
\tikzset{dot/.style={circle,fill,inner sep=1.5pt},
pics/momentum arrow/.style={code={%
\def\pv##1{\pgfkeysvalueof{/tikz/momentum arrow/##1}}
\tikzset{/tikz/momentum arrow/.cd,#1}
\pgfmathtruncatemacro{\itest}{\pv{d}>0}
\ifnum\itest>0
\draw[momentum arrow/marrow] (-0.5*\pv{l},\pv{d}) -- (0.5*\pv{l},\pv{d})
node[midway,auto]{\pv{m}};
\else
\draw[momentum arrow/marrow] (-0.5*\pv{l},\pv{d}) -- (0.5*\pv{l},\pv{d})
node[midway,auto]{\pv{m}};
\fi
}},
pics/momentum arrow'/.style={code={%
\def\pv##1{\pgfkeysvalueof{/tikz/momentum arrow/##1}}
\tikzset{/tikz/momentum arrow/.cd,#1}
\pgfmathtruncatemacro{\itest}{\pv{d}>0}
\ifnum\itest>0
\draw[momentum arrow/marrow'] (-0.5*\pv{l},\pv{d}) -- (0.5*\pv{l},\pv{d})
node[midway,auto]{\pv{m}};
\else
\draw[momentum arrow/marrow'] (-0.5*\pv{l},\pv{d}) -- (0.5*\pv{l},\pv{d})
node[midway,auto]{\pv{m}};
\fi
}},momentum arrow/.cd,d/.initial=5mm,l/.initial=1cm,
marrow/.style={-latex},marrow'/.style={latex-},m/.initial={}}
\tikzset{pics/fermion flow/.style={code={%
\def\pv##1{\pgfkeysvalueof{/tikz/fermion flow/##1}}
\tikzset{/tikz/fermion flow/.cd,#1}
\pgfmathtruncatemacro{\itest}{\pv{d}>0}
\ifnum\itest>0
\draw[fermion flow/fflow] (-0.5*\pv{l},\pv{d}) -- (0.5*\pv{l},\pv{d})
node[midway,auto]{\pv{m}};
\else
\draw[fermion flow/fflow] (-0.5*\pv{l},\pv{d}) -- (0.5*\pv{l},\pv{d})
node[midway,auto]{\pv{m}};
\fi
}},fermion flow/.cd,d/.initial=-5mm,l/.initial=1cm,
fflow/.style={-{Latex[bend]},gray!50},m/.initial={}}
% \pgfdeclaredecoration{complete sines}{initial}
% {
%     \state{initial}[
%         width=+0pt,
%         next state=sine,
%         persistent precomputation={\pgfmathsetmacro\matchinglength{
%             \pgfdecoratedinputsegmentlength / int(\pgfdecoratedinputsegmentlength/\pgfdecorationsegmentlength)}
%             \setlength{\pgfdecorationsegmentlength}{\matchinglength pt}
%         }] {}
%     \state{sine}[width=\pgfdecorationsegmentlength]{
%         \pgfpathsine{\pgfpoint{0.25\pgfdecorationsegmentlength}{0.5\pgfdecorationsegmentamplitude}}
%         \pgfpathcosine{\pgfpoint{0.25\pgfdecorationsegmentlength}{-0.5\pgfdecorationsegmentamplitude}}
%         \pgfpathsine{\pgfpoint{0.25\pgfdecorationsegmentlength}{-0.5\pgfdecorationsegmentamplitude}}
%         \pgfpathcosine{\pgfpoint{0.25\pgfdecorationsegmentlength}{0.5\pgfdecorationsegmentamplitude}}
% }
%     \state{final}{}
% }


\endinput
